% !TeX root = ../动手学强化学习笔记.tex
% 由主文档通过 \input 载入；请编译主文档。
\setcounter{chapter}{0} % 保留原书章号。
\chapter{绪论}
\label{ch:rl-introduction}

\textbf{学习目标与本章安排。}从可手算的决策例子认识强化学习的研究对象，区分即时奖励与长期回报、环境状态与观测、策略与学习算法。本章不依赖后续算法；先建立直观认识，再整理基本概念，最后给出问题分类与学习检查。

强化学习（Reinforcement Learning，RL）是一类研究\emph{智能体如何通过与环境交互来学习决策策略}的机器学习方法。在交互过程中，智能体观察环境当前的状态，选择并执行一个动作，随后环境转移到新的状态，并向智能体反馈一个数值奖励。智能体根据不断积累的交互经验调整自己的行为，其目标不是只追求某一步的即时奖励，而是学习一个能够使\emph{长期累积奖励最大化}的策略。

\section{从两条路线理解长期决策}

\begin{example}[先付出代价是否值得]
在一个确定性任务中，智能体从起点选择两条路线之一。路线甲立刻获得奖励 $2$，随后任务终止；路线乙第一步获得奖励 $-1$，第二步必然获得奖励 $5$，然后终止。若折扣因子为 $\gamma=0.9$，两条路线的回报分别为 $2$ 和 $-1+0.9\times5=3.5$，因而路线乙的长期收益更高。若 $\gamma=0.5$，路线乙的回报变成 $1.5$，最优选择随目标的改变而改变。这个例子表明，决策是否合理必须相对于完整任务及其回报定义判断。
\end{example}

在这个例子中，用状态、动作、转移和奖励描述两条路线，是建立数学模型；编写程序让动作产生相应的下一状态与奖励，是实现环境；利用多次试验改进路线选择，是学习算法的工作；训练后在起点选择哪条路线，则由策略决定。若模型完全已知，可以直接计算路线回报；若结果未知，才需要通过数据估计。这个区别将在后续各章反复使用。


\section{强化学习的基本概念}


强化学习研究的是\emph{不确定环境中的序贯决策问题}。与针对每个输入直接给出正确标签的监督学习不同，强化学习中的智能体通常只接收标量形式的奖励信号，并且当前动作不仅影响即时奖励，还会改变未来所处的状态以及此后能够获得的全部奖励。因此，强化学习问题具有三个相互关联的核心特征：\emph{反馈具有评价性而非指示性、动作具有延迟影响、数据分布会随策略改变}。智能体必须在利用已有知识与探索未知行为之间取得平衡，并从交互产生的经验中改进决策规则。


\subsection{智能体与环境的交互过程}


强化学习通常在离散决策时刻 $t=0,1,2,\ldots$ 上描述。时刻 $t$，环境具有状态 $S_t$，智能体接收到观测 $O_t$，并依据其策略选择动作 $A_t$。环境执行该动作后转移到状态 $S_{t+1}$，同时产生奖励 $R_{t+1}$；随后智能体接收到新的观测 $O_{t+1}$。这一交互闭环可以写成
\[
  O_t \longrightarrow A_t \longrightarrow (R_{t+1},O_{t+1})
  \longrightarrow A_{t+1} \longrightarrow \cdots .
\]
若显式保留环境状态，则一次交互生成的轨迹可表示为
\[
  \tau=(S_0,O_0,A_0,R_1,S_1,O_1,A_1,R_2,\ldots).
\]
这里采用强化学习文献中的常见时间约定：$A_t$ 是智能体在时刻 $t$ 采取的动作，$R_{t+1}$ 是环境对该动作给出的下一时刻奖励。随机变量使用大写字母表示，其具体取值使用相应的小写字母表示，例如 $S_t=s$、$A_t=a$。


在一般的部分可观测情形下，智能体在时刻 $t$ 可利用的信息历史为
\[
  H_t=(O_0,A_0,R_1,O_1,\ldots,A_{t-1},R_t,O_t).
\]
智能体按照策略从 $H_t$ 选择动作，环境则根据其内部状态和动作产生下一状态与奖励。一个常用的概率描述是
\begin{align*}
  A_t &\sim \pi(\,\cdot\mid H_t),\\
  (S_{t+1},R_{t+1}) &\sim p(\,\cdot,\cdot\mid S_t,A_t),\\
  O_{t+1} &\sim \Omega(\,\cdot\mid S_{t+1}),
\end{align*}
其中 $p$ 是环境的状态转移与奖励核，$\Omega$ 是观测核，$\pi$ 是智能体的策略。这三个概率核分别刻画环境如何演化、智能体能够看到什么以及智能体如何行动。


\subsection{七个基本要素}


\begin{enumerate}[label=\textbf{\arabic*.},leftmargin=3em]
  \item \textbf{智能体（agent）}


  智能体是进行决策与学习的主体。它接收来自环境的信息，选择动作，并根据交互经验更新自身的决策机制。智能体是一个由建模边界所确定的概念，不必等同于完整的物理实体：在机器人控制中，机器人本体、传感器和执行器通常属于环境的一部分，而控制算法可被视为智能体；在游戏中，智能体可以是操纵角色的程序；在自动驾驶中，智能体可以是负责规划与控制的决策系统。


  从算法角度看，智能体可以包含策略、价值函数、环境模型、记忆状态以及探索机制等组件，但只有将观测或历史映射为动作的决策机制是定义强化学习智能体所必需的。


  \item \textbf{环境（environment）}


  环境是位于智能体决策边界之外、能够接收智能体动作并对其作出响应的系统。它决定状态如何演化、智能体获得何种观测，以及产生怎样的奖励。环境既可以是物理世界，也可以是仿真器、博弈规则、市场过程或数据生成机制。


  在完全可观测的马尔可夫模型中，环境的单步动力学常由联合概率核
  \[
    p(s',r\mid s,a)
    =\Pr\!\left(S_{t+1}=s',R_{t+1}=r\mid S_t=s,A_t=a\right)
  \]
  描述。对于连续状态或连续奖励，上式应理解为相应的概率密度或一般概率核。环境可以是确定性的，也可以是随机的；智能体通常并不知道 $p$ 的具体形式，而需要通过经验学习如何决策。


  \item \textbf{状态（state，$s$）}


  状态是对环境在某一时刻所处情形的表示，状态随机变量记为 $S_t$，其取值 $s$ 属于状态空间 $\mathcal{S}$。在强化学习理论中，理想状态应当包含预测未来所需的全部相关信息。若 $S_t$ 满足马尔可夫性质，则在给定当前状态与动作后，未来与更早的历史条件独立：
  \[
    \Pr(S_{t+1},R_{t+1}\mid S_0,A_0,R_1,\ldots,S_t,A_t)
    =\Pr(S_{t+1},R_{t+1}\mid S_t,A_t).
  \]
  因而，马尔可夫状态是关于未来动力学的\emph{充分信息表示}，而不只是某些传感器读数的简单集合。需要区分环境的真实状态、研究者为建模引入的状态，以及智能体内部维护的信息状态；三者并不必然相同。


  \item \textbf{观测（observation，$o$）}


  观测是智能体在某一时刻实际获得的信息，观测随机变量记为 $O_t$，其取值 $o$ 属于观测空间 $\mathcal{O}$。观测可以由状态经过观测核 $\Omega(o\mid s)$ 随机生成，也可能包含测量噪声、信息缺失或感知误差。


  在完全可观测的马尔可夫决策过程中，智能体能够获得足以确定马尔可夫状态的信息，通常简写为 $O_t=S_t$。在部分可观测的情形下，单个观测不足以预测未来，问题应建模为部分可观测马尔可夫决策过程。此时策略一般需要依赖历史 $H_t$，或者依赖关于隐状态的信念分布
  \[
    b_t(s)=\Pr(S_t=s\mid H_t).
  \]
  因此，\emph{状态是用于描述环境和预测未来的变量，观测是智能体真正能够访问的信息}，二者不能在一般情形下混为一谈。


  \item \textbf{动作（action，$a$）}


  动作是智能体能够施加于环境的决策或控制输入。动作随机变量记为 $A_t$，其取值 $a$ 属于动作空间 $\mathcal{A}$；若可行动作依赖于当前状态，则可写为 $a\in\mathcal{A}(s)$。动作空间可以是离散的，例如棋类游戏中的合法落子；也可以是连续的，例如机器人关节力矩；还可以是同时包含离散选择和连续参数的混合空间。


  动作会同时影响即时奖励和未来状态分布，这正是强化学习区别于静态预测问题的关键。一个当前看来收益较低的动作，可能通过到达更有利的状态而提高长期收益。


  \item \textbf{奖励（reward，$r$）}


  奖励是环境在一次状态转移后向智能体提供的标量反馈，奖励随机变量通常记为 $R_{t+1}$。它刻画任务设计者希望智能体优化的目标，但只评价当前转移产生的即时结果。常用的期望奖励函数为
  \[
    r(s,a)
    \coloneqq \mathbb{E}\!\left[R_{t+1}\mid S_t=s,A_t=a\right],
  \]
  当奖励还显式依赖下一状态时，也可写为 $r(s,a,s')$。奖励可能稠密或稀疏、确定或随机、正或负；负奖励常被解释为代价，但奖励的正负本身没有独立于任务目标的道德含义。


  必须严格区分\emph{奖励}与\emph{回报}：奖励 $R_{t+1}$ 是单步反馈，回报则是从某个时刻起累积的未来奖励。智能体的目标不是逐步贪心地最大化每一个即时奖励，而是最大化期望长期回报。


  \item \textbf{策略（policy，$\pi$）}


  策略是智能体的行为规则，它规定在给定可用信息时如何选择动作。对于完全可观测且采用平稳马尔可夫策略的离散问题，随机策略定义为条件概率分布
  \[
    \pi(a\mid s)
    \coloneqq \Pr(A_t=a\mid S_t=s),
    \qquad
    \sum_{a\in\mathcal{A}(s)}\pi(a\mid s)=1.
  \]
  确定性策略可写为映射
  \[
    A_t=\mu(S_t),
    \qquad \mu:\mathcal{S}\to\mathcal{A},
  \]
  它为每个状态直接指定一个动作；随机策略则为可行动作给出概率分布。在部分可观测问题中，更一般的策略应写为 $\pi(a\mid H_t)$，或者写为依赖当前观测、记忆状态或信念状态的形式。实际算法中，策略往往由参数 $\theta$ 表示为 $\pi_\theta$，学习的任务就是利用交互数据调整 $\theta$，使长期性能指标得到改善。
\end{enumerate}


\subsection{回报与强化学习目标}


为了形式化“长期累积奖励”，定义时刻 $t$ 的折扣回报为
\begin{equation}
  G_t
  \coloneqq
  \sum_{k=0}^{\infty}\gamma^k R_{t+k+1},
  \qquad 0\leq \gamma<1,
  \label{eq:discounted-return}
\end{equation}
其中 $\gamma$ 称为折扣因子。若奖励几乎处处满足 $|R_{t+1}|\leq R_{\max}<\infty$，则 $|G_t|\leq R_{\max}/(1-\gamma)$，从而该级数绝对收敛且可积。较大的 $\gamma$ 使智能体更加重视长期后果；较小的 $\gamma$ 则使近期奖励占据更大权重。在有限时域的情节式任务中，求和在终止时刻结束，并且在回报可积时可以取 $\gamma=1$。


给定初始状态分布以及环境动力学，策略 $\pi$ 的典型优化目标为
\begin{equation}
  J(\pi)
  \coloneqq
  \mathbb{E}_{\pi}\!\left[G_0\right],
  \label{eq:rl-objective}
\end{equation}
其中下标 $\pi$ 表明期望同时关于策略引入的随机性、环境转移、奖励、观测以及初始状态取值。于是，强化学习可以概括为：\emph{智能体在未知或不完全已知的环境中，通过交互数据学习策略，使某种期望累积回报指标最大化。}对于无限时域持续任务，还可以采用平均奖励等其他目标；选择何种目标属于问题建模的一部分。


\begin{remark}
奖励函数规定“希望实现什么”，策略规定“智能体如何行动”，环境动力学规定“采取动作之后会发生什么”。三者承担不同角色。一个形式正确的强化学习问题必须明确智能体与环境的边界、智能体可获得的信息、允许采取的动作、奖励产生机制以及需要优化的长期性能指标。
\end{remark}


\section{强化学习问题的分类与方法地图}

\emph{预测}（prediction）是在策略固定时估计未来回报；\emph{控制}（control）是寻找表现更好的策略。以迷宫为例，估计“始终沿右侧墙壁走”需要多少步属于预测，寻找更快到达出口的规则属于控制。模型已知时，可通过计算比较动作后果，这称为规划（planning）；模型未知时，则必须从交互经验获得有关后果的信息。学习与规划可以结合，而不是两个互斥的任务。

\emph{基于模型方法}（model-based method）利用环境模型预测动作后果并支持决策；\emph{无模型方法}（model-free method）直接从经验更新价值或策略，不显式学习用于规划的转移模型。二者都需要环境信息。\emph{同策略学习}（on-policy learning）评价或改进产生当前数据的策略；\emph{离策略学习}（off-policy learning）允许目标策略不同于行为策略。\emph{在线学习}（online learning）可以在学习过程中继续收集交互数据；\emph{离线学习}（offline learning）只使用给定数据集。这些分类回答不同问题：离策略算法既可以在线运行，也可以用于离线数据，但后者需要额外处理覆盖不足。

表格方法为每个状态或状态动作对保存独立数值，适合小规模问题；函数近似用共享参数表示价值或策略，能够在状态之间泛化。神经网络只是函数近似的一种形式，并不改变强化学习所优化的长期决策目标。相较于通常使用固定样本的监督学习，强化学习中的动作会改变后续数据分布，因此数据获取与模型优化相互影响。

\section{典型困难与建模检查}

\emph{探索与利用}（exploration and exploitation）分别强调获取未知动作的信息与利用现有估计获得收益。当前估计最好的动作未必真正最好，但探索也会消耗预算。\emph{时间信用分配}（temporal credit assignment）回答成功或失败应归因于哪些先前动作；终点的一次奖励通常无法直接告诉智能体哪一步最关键。稀疏奖励进一步使大量不同轨迹具有相同反馈，而部分可观测性要求智能体利用记忆辨认当前处境。

一个可检验的任务定义必须明确可用信息、动作限制、奖励产生方式、初始状态分布、终止规则及评价目标。若机器人因传感器缺失看不到障碍物，不能在策略输入中偷偷加入完整地图；若训练按折扣回报优化而评价按成功率排名，两者的差异必须说明。奖励是对任务要求的数学编码，最大化奖励只能保证与该编码一致，不能保证编码本身完整表达了设计者的意图。

\section{习题}
\begin{enumerate}
\item 在两条路线的例子中求使两种动作同样好的折扣因子，分别给出其两侧的最优动作。若第二条路线第二步只有概率 $0.8$ 获得 $5$，否则获得 $0$，阈值如何改变？
\item 为带有限电量的扫地机器人给出状态、观测、动作、奖励和终止条件。说明仅用当前位置作为状态可能遗漏什么。
\item 分别说明“MDP 是学习算法”“奖励就是正确动作标签”“无模型方法不需要环境数据”错在哪里，并为在线离策略学习给出一个具体例子。
\end{enumerate}





